% Chapter 6: Conclusion and Future Work

\chapter{Conclusion and Future Work}
\label{Chapter6}

\section{Conclusion}
This thesis designed, implemented, and evaluated a real-time, browser-accessible Indian Sign Language recognition and bidirectional translation system. By extracting 21 3D landmarks per hand and engineering a 156-dimensional scale- and translation-invariant feature representation, the platform captures the bimanual structures characteristic of ISL.

The ensemble Random Forest classifier achieves 85.2\% held-out test accuracy across 45 gesture classes with 1.2 ms inference latency and 27.2 ms end-to-end processing time ($\approx 36.8$ FPS). The temporal hysteresis engine eliminates prediction flicker, while the reverse translation pathway combines a 204-entry sign dictionary, 203 demonstration videos, speech recognition, and a Deepgram Nova-2 YouTube pipeline.

\section{Practical Applications}
The platform is suited for:
\begin{itemize}
\item Public administrative desks, banking counters, and hospital intake stations.
\item Self-paced ISL educational practice using split-screen camera feedback.
\item Classroom and workplace accessibility support.
\end{itemize}

\section{Limitations and Future Work}
Future research directions include:
\begin{enumerate}
\item \textbf{Continuous Sign Grammar:} Transitioning from isolated gestures to continuous sentences using spatio-temporal Transformers and Connectionist Temporal Classification (CTC).
\item \textbf{Holistic Multi-Modal Cues:} Incorporating facial expressions, mouth shapes, and upper-body pose via MediaPipe Holistic.
\item \textbf{Client-Side Acceleration:} Porting inference to TensorFlow.js, ONNX Runtime Web, and WebGPU for zero-server-latency execution.
\item \textbf{Community Validation:} Conducting formal user studies with fluent deaf ISL signers to refine linguistic fidelity.
\end{enumerate}
